\documentclass[10pt,letterpaper]{article}

\usepackage[T1]{fontenc}
\usepackage[utf8]{inputenc}
\usepackage{lmodern}
\usepackage[margin=0.85in]{geometry}
\usepackage{enumitem}
\usepackage{hyperref}
\usepackage{fancyhdr}
\usepackage{lastpage}
\usepackage{xcolor}

\definecolor{darknavy}{HTML}{0000A3}
\hypersetup{colorlinks=true, linkcolor=darknavy, urlcolor=darknavy, citecolor=darknavy}
\pagestyle{fancy}
\fancyhf{}
\renewcommand{\headrulewidth}{0pt}
\renewcommand{\footrulewidth}{0pt}
\fancyfoot[R]{\small \thepage}

\newcommand{\sectionheader}[1]{\vspace{0.55em}\noindent\textbf{#1}\vspace{0.15em}\par}
\newcommand{\pubitem}[2]{\item[#1] #2}

\begin{document}
\thispagestyle{empty}

\noindent{\Large\textbf{Dar\'io Gonz\'alez Picos}}\\[0.15em]
\noindent PhD Candidate, Leiden Observatory, Leiden University, The Netherlands\\
\noindent \href{mailto:picos@strw.leidenuniv.nl}{picos@strw.leidenuniv.nl} ---
\href{https://dgonzalezpicos.github.io}{dgonzalezpicos.github.io} ---
\href{https://github.com/DGonzalezPicos}{github.com/DGonzalezPicos}

\vspace{0.6em}
\noindent PhD candidate in exoplanet atmospheres, focusing on high-resolution spectroscopy and atmospheric characterization of exoplanets and host stars.\\
Experienced in observational astronomy, data reduction, and spectral analysis of medium- to high-resolution spectroscopic data from ground-based telescopes and JWST.\\
As of 2026-06-11, part of 15 refereed papers with 206 total citations; a complete listing can be found in my \href{https://ui.adsabs.harvard.edu/search/p_=0&q=author%3A%22Gonz%C3%A1lez%20Picos%2C%20D.%3B%20Picos%2C%20Dar%C3%ADo%20Gonz%C3%A1lez%22&sort=date%20desc%2C%20bibcode%20desc}{ADS} or \href{https://orcid.org/0000-0001-9282-9462}{ORCID}.

\sectionheader{Experience}
\noindent\textbf{Leiden Observatory, Leiden University, The Netherlands}\\
\noindent 2023 -- present \hspace{1.5em} PhD Candidate\\
Researching exoplanet atmospheres using high-resolution spectroscopy, focusing on linking formation pathways and atmospheric composition of Super-Jupiters via the CO isotopologue ratio. Advisor: Prof.\ Dr.\ I.A.G Snellen.

\vspace{0.35em}
\noindent\textbf{Leiden University, The Netherlands}\\
\noindent Feb 2022 -- present \hspace{0.5em} Student Teaching Assistant\\
Tutored problem sessions of Master's course Star and Planet Formation, Bachelor's course Radiative Processes in Astrophysics.

\vspace{0.35em}
\noindent\textbf{Isaac Newton Group of Telescopes (ING), La Palma, Spain}\\
\noindent Sep 2020 -- Sep 2021 \hspace{0.5em} Student Support Astronomer\\
Support Astronomer and Telescope Operator at the 2.5-m INT at Observatorio Roque de Los Muchachos (+40 observing nights).

\sectionheader{Education}
\noindent PhD 2023--2027 (expected), Leiden Observatory, Leiden University, The Netherlands. Advisor: Snellen\\
MSc 2021--2023, Astronomy, Leiden University, The Netherlands\\
BSc 2016--2020, Physics, University of Barcelona, Spain

\sectionheader{Selected Publications}
\begin{enumerate}[label=\arabic*., leftmargin=2.2em, itemsep=0.25em, topsep=0.15em, parsep=0pt]
\pubitem{7}{\textbf{Gonz\'alez Picos, D.}; Snellen, I. A. G.; Landman, R.; de Regt, S.; Grasser, N.; et al., 2026, \href{https://arxiv.org/abs/2606.12401}{The carbon isotope ratio of $\beta$ Pic b with high-resolution spectroscopy}, \emph{Astronomy and Astrophysics} (accepted; arXiv:2606.12401)}
\pubitem{6}{\textbf{Gonz\'alez Picos, D.}; van der Post, T.; de Regt, S.; Ruffio, J.-B.; Grasser, N.; \& Snellen, I.A.G., 2026, \href{https://doi.org/10.1051/0004-6361/202659296}{JWST high-contrast spectroscopy with speckle modelling: Atmospheric retrievals of the T-dwarf companion HD 19467 B}, \emph{Astronomy and Astrophysics} (arXiv:2605.02769)}
\pubitem{5}{\textbf{Gonz\'alez Picos, Dar\'io}; Snellen, Ignas; \& de Regt, Sam, 2025, \href{https://doi.org/10.1038/s41550-025-02641-4}{Chemical evolution imprints in the rare isotopes of nearby M dwarfs}, \emph{Nature Astronomy}, 9, 1692 (arXiv:2508.18424) [5 citations]}
\pubitem{4}{\textbf{Gonz\'alez Picos, D.}; de Regt, S.; Gandhi, S.; Grasser, N.; \& Snellen, I. A. G., 2025, \href{https://doi.org/10.1051/0004-6361/202555669}{Disentangling disc and atmospheric signatures of young brown dwarfs with JWST/NIRSpec}, \emph{Astronomy and Astrophysics}, 703 (arXiv:2509.11694)}
\pubitem{3}{\textbf{Gonz\'alez Picos, D.}; Snellen, I. A. G.; de Regt, S.; Landman, R.; et al., 2025, \href{https://doi.org/10.1051/0004-6361/202451936}{The ESO SupJup Survey: IV. Unveiling the carbon isotope ratio of GQ Lup B and its host star}, \emph{Astronomy and Astrophysics}, 693 (arXiv:2501.01789) [16 citations]}
\pubitem{2}{Zhang, Yapeng; \textbf{Gonz\'alez Picos, Dar\'io}; de Regt, Sam; Snellen, Ignas A. G.; et al., 2024, \href{https://doi.org/10.3847/1538-3881/ad7ea9}{The ESO SupJup Survey. III. Confirmation of $^{13}$CO in YSES 1 b and Atmospheric Detection of YSES 1 c with CRIRES$^{+}$}, \emph{The Astronomical Journal}, 168, 246 (arXiv:2409.16660) [33 citations]}
\pubitem{1}{\textbf{Gonz\'alez Picos, D.}; Snellen, I. A. G.; de Regt, S.; Landman, R.; et al., 2024, \href{https://doi.org/10.1051/0004-6361/202450028}{The ESO SupJup Survey: II. The $^{12}$C/$^{13}$C isotope ratios of three young brown dwarfs with CRIRES$^{+}$}, \emph{Astronomy and Astrophysics}, 689 (arXiv:2407.07678) [22 citations]}
\end{enumerate}

\sectionheader{Talks}
\noindent Invited talk --- CNOPS26, Sexten, Italy: C and O isotopic ratios in planets and host stars. [\href{https://dgonzalezpicos.github.io}{slides}]\\
Contributed talk --- AAS247, Phoenix, USA: Chemical evolution imprints in the rare isotopes of nearby M dwarfs. [\href{https://ui.adsabs.harvard.edu/abs/2026AAS...24723101G/abstract}{abstract}]\\
Contributed talk --- ExoELT2025, Garching ESO, Germany: Storm chasers: Probing $\beta$ Pic b's weather with high-resolution spectroscopy.\\
Flash talk --- EAS 2025, Cork, Ireland: Rare isotopes in nearby M dwarfs probe galactic chemical evolution.\\
Contributed talk --- ESA Holland-Area Exoplanet Science Meeting, ESA/ESTEC, The Netherlands: JWST/NIRSpec's full 0.97--5.27\,$\mu$m view of young brown dwarfs.\\
Contributed talk --- STRW Science Day, Leiden Observatory, The Netherlands: Carbon isotope ratios in exoplanet and brown dwarf atmospheres.\\
Flash talk --- Two HorSeS, Berlin, Germany: Atmospheric Characterization and Isotopic Measurements of Directly-Imaged Planets and Brown Dwarfs with VLT/CRIRES$^{+}$.\\
Flash talk --- Thinkshop 2022 HRS for Exoplanet Atmospheres, Potsdam, Germany: Cross-correlation of exoplanet atmospheres at high-resolution.

\vfill
\begin{center}
Last updated: June 11, 2026
\end{center}

\end{document}
